\documentclass[11pt,a4paper]{article}
\usepackage[margin=1in]{geometry}
\usepackage{amsmath,amssymb,amsthm}
\usepackage{algorithm}
\usepackage{algpseudocode}
\usepackage{booktabs}
\usepackage{hyperref}

\theoremstyle{definition}
\newtheorem{definition}{\textbf{Definition}}[section]
\newtheorem{theorem}{\textbf{Theorem}}[section]
\newtheorem{lemma}[theorem]{\textbf{Lemma}}
\newtheorem{remark}{\textbf{Remark}}[section]

\title{\textbf{ZeRO and Fully Sharded Data Parallelism (FSDP): Linear Model State Partitioning}}
\author{\textbf{Team Scaibu} \\ \small Email: \texttt{jaydeep.9664920749@gmail.com} \\ \small \textbf{Architecture and Core Engine Team}}
\date{\textbf{October 2026}}

\begin{document}
\maketitle

\section{Definitions and Cognitive Mental Models}

\subsection{Formal Mathematical Definition}
\textbf{Definition (Zero Redundancy Optimizer and Fully Sharded Partitioning):}
Standard Data Parallelism (DDP) replicates the entire model state—parameters $\theta \in \mathbb{R}^P$, gradients $g \in \mathbb{R}^P$, and optimizer states (for Adam: FP32 master weights $W_{\text{FP32}}$, first moment $m$, and second moment $v$)—across every GPU device. For a model with $P$ parameters trained in 16-bit precision, static model states consume $2P + 2P + 12P = 16P$ bytes per device, creating a fundamental memory barrier for models exceeding several billion parameters.

\textbf{ZeRO (Zero Redundancy Optimizer)} (Rajbhandari et al., 2020) and \textbf{Fully Sharded Data Parallelism (FSDP)} eliminate this replication by sharding model states across the $N_{\text{dev}}$ data-parallel processes:
\begin{enumerate}
    \item \textbf{ZeRO-Stage 1 ($P_{\text{os}}$)}: Shards optimizer states across $N_{\text{dev}}$ devices. Each device stores only $12P / N_{\text{dev}}$ bytes of optimizer state while replicating parameters ($2P$) and gradients ($2P$).
    \item \textbf{ZeRO-Stage 2 ($P_{\text{os+g}}$)}: Shards optimizer states and gradients across $N_{\text{dev}}$ devices. During backpropagation, gradients are reduce-scattered, so each device retains only $2P / N_{\text{dev}}$ gradient bytes.
    \item \textbf{ZeRO-Stage 3 / FSDP ($P_{\text{os+g+p}}$)}: Shards optimizer states, gradients, and parameters. Each device stores only $2P / N_{\text{dev}}$ parameter bytes.
\end{enumerate}
Under ZeRO-Stage 3, parameters for layer $l$ are transiently reconstructed on-the-fly via an \textbf{All-Gather} collective immediately before layer $l$'s forward pass, freed from memory upon forward completion, and re-gathered during the backward pass. Per-device memory footprint scales as:
{\boldmath
\begin{equation}
\text{Memory}_{\text{ZeRO-3}} = \frac{16 P}{N_{\text{dev}}} \cdot B_{\text{bytes}}
\end{equation}
}
achieving a linear $N_{\text{dev}}\times$ reduction in memory requirements.

\subsection{Expanded Non-Technical Definition (Intuition for Anyone)}
\begin{enumerate}
    \item \textbf{Shared Book, Not 8 Copies}: In standard training, if you have 8 GPUs, every single GPU holds an identical 100-page book in its small backpack. ZeRO-3 tears the book into 8 pieces: GPU 1 carries pages 1–12, GPU 2 carries pages 13–25, and so on.
    \item \textbf{Pass It When You Read It}: When GPU 1 reaches chapter 2, all GPUs quickly swap photocopies of chapter 2, do their calculations, and immediately throw the photocopies away into the recycling bin!
    \item \textbf{Train Giant Models on Ordinary Clusters}: Without ZeRO, a 70-billion parameter model requires an expensive supercomputer with massive GPU memory. With ZeRO-3/FSDP, that same model can be sharded across standard GPUs because memory scales down with the number of GPUs.
\end{enumerate}

\subsection{Child-Friendly Mental Model (Physical Metaphor)}
Imagine 8 students doing a massive jigsaw puzzle:
\begin{itemize}
    \item \textbf{Standard DDP}: Buying 8 complete identical puzzle boxes. Every student's desk is overflowing with all 10,000 pieces!
    \item \textbf{ZeRO-3 / FSDP}: Buying one puzzle box and giving 1,250 pieces to each student. When student 1 needs a corner piece, they borrow it for 5 seconds and give it right back.
\end{itemize}

\section{Core Mathematical Equations and Definitions}

\subsection{ZeRO Memory Partitioning Hierarchy}
{\boldmath
\begin{align}
\mathcal{M}_{\text{DDP}} &= 2P + 2P + 12P = 16P \text{ bytes} \\
\mathcal{M}_{\text{ZeRO-1}} &= 2P + 2P + \frac{12P}{N} = 4P + \frac{12P}{N} \text{ bytes} \\
\mathcal{M}_{\text{ZeRO-2}} &= 2P + \frac{2P}{N} + \frac{12P}{N} = 2P + \frac{14P}{N} \text{ bytes} \\
\mathcal{M}_{\text{ZeRO-3}} &= \frac{2P}{N} + \frac{2P}{N} + \frac{12P}{N} = \frac{16P}{N} \text{ bytes}
\end{align}
}
\begin{itemize}
    \item \textbf{Formal Definition}: Memory scaling equations for baseline DDP and ZeRO Stages 1, 2, and 3.
    \item \textbf{What It Means}: Quantifies exact memory capacity requirements per accelerator.
    \item \textbf{Why It Is Important}: Directly guides cluster sizing and parallelism layout design.
\end{itemize}

\subsection{Communication Volume Expansion Factor}
{\boldmath
\begin{equation}
\text{Comm}_{\text{ZeRO-3}} = 1.5 \times \text{Comm}_{\text{DDP}}
\end{equation}
}
\begin{itemize}
    \item \textbf{Formal Definition}: Communication overhead of ZeRO-Stage 3 relative to standard data-parallel all-reduce.
    \item \textbf{What It Means}: ZeRO-3 adds an all-gather in the forward pass, increasing communication by exactly 50\%.
    \item \textbf{Why It Is Important}: Explains why prefetching is critical to hide communication latency behind compute.
\end{itemize}

\section{Rigorous Mathematical Analysis Checklist}

\begin{itemize}
    \item \textbf{Notation}: $P$ (total parameters), $N$ (device ranks), $B$ (bytes per parameter).
    \item \textbf{Epistemic Status}: Rajbhandari et al. (SC 2020 - DeepSpeed); Zhao et al. (PyTorch FSDP). Foundation of modern distributed LLM training.
    \item \textbf{Dependency Graph}: Sharded Weights $\to \text{Prefetch All-Gather} \to \text{Forward Layer } l \to \text{Free Weights} \to \text{Backward Layer } l \to \text{Reduce-Scatter Grads}$.
    \item \textbf{Dimensions}: Model size $P \in [10^9, 10^{12}]$ parameters, cluster ranks $N \in [8, 16384]$.
    \item \textbf{Regimes}: Large-scale pre-training and full fine-tuning of 7B to 405B models.
    \item \textbf{Invariants}: At step completion, every device's local optimizer step produces mathematically identical weight updates to monolithic training.
    \item \textbf{Isomorphisms}: Isomorphic to partitioned address space shared-memory multiprocessor cache hierarchies.
\end{itemize}

\section{Theoretical Theorems and Formal Proofs}

\begin{theorem}[Linear Memory Scaling of ZeRO-Stage 3]
Let model parameter count be $P$ and precision byte size be $B = 2$ bytes. Under ZeRO-Stage 3 parameter, gradient, and optimizer state sharding, the static memory per device satisfies:
{\boldmath
\begin{equation}
\lim_{N \to \infty} \mathcal{M}_{\text{ZeRO-3}}(N) = 0
\end{equation}
}
and for any target memory budget $\mathcal{M}_{\text{target}}$, there exists a finite device count $N^* = \lceil 16P / \mathcal{M}_{\text{target}} \rceil$ such that the model fits in GPU memory.
\end{theorem}

\begin{proof}
By definition of ZeRO-Stage 3:
- Parameters are partitioned into $N$ disjoint shards of size $P/N$. At 16-bit precision, this consumes $2P/N$ bytes.
- Gradients are accumulated locally and immediately reduce-scattered into $N$ shards of size $P/N$, consuming $2P/N$ bytes.
- Adam optimizer states (FP32 master weight 4B + momentum 4B + variance 4B = 12B/param) are stored exclusively for the local $P/N$ parameter slice, consuming $12P/N$ bytes.
Summing all static components:
{\boldmath
\begin{equation}
\mathcal{M}_{\text{static}}(N) = \frac{2P}{N} + \frac{2P}{N} + \frac{12P}{N} = \frac{16P}{N}
\end{equation}
}
Taking the limit as $N \to \infty$:
{\boldmath
\begin{equation}
\lim_{N \to \infty} \frac{16P}{N} = 0
\end{equation}
}
Setting $\frac{16P}{N} \le \mathcal{M}_{\text{target}}$ and solving for $N$ yields:
{\boldmath
\begin{equation}
N \ge \frac{16P}{\mathcal{M}_{\text{target}}} \implies N^* = \left\lceil \frac{16P}{\mathcal{M}_{\text{target}}} \right\rceil
\end{equation}
}
Thus, any arbitrarily large model can be trained on fixed-memory devices simply by scaling the device count $N$.
\end{proof}

\begin{theorem}[Exact $1.5\times$ Communication Multiplier of ZeRO-3]
The total network communication volume per step in ZeRO-Stage 3 is exactly $1.5\times$ that of standard Data Parallelism with Ring All-Reduce:
{\boldmath
\begin{equation}
\frac{\text{Comm}_{\text{ZeRO-3}}}{\text{Comm}_{\text{DDP}}} = 1.5
\end{equation}
}
\end{theorem}

\begin{proof}
Let total parameter size be $M$ bytes.
1. In standard DDP:
Gradients of size $M$ are aggregated using Ring All-Reduce:
- Reduce-Scatter phase transmits $\frac{N-1}{N} M$ bytes.
- All-Gather phase transmits $\frac{N-1}{N} M$ bytes.
Total DDP communication per step:
{\boldmath
\begin{equation}
\text{Comm}_{\text{DDP}} = 2 \left( \frac{N-1}{N} \right) M
\end{equation}
}
2. In ZeRO-Stage 3:
- In the forward pass, layer parameters are reconstructed via All-Gather: $\left( \frac{N-1}{N} \right) M$ bytes.
- In the backward pass, layer parameters are re-gathered via All-Gather: $\left( \frac{N-1}{N} \right) M$ bytes.
- In the backward pass, gradients are reduce-scattered to shard owners: $\left( \frac{N-1}{N} \right) M$ bytes.
Summing all three phases:
{\boldmath
\begin{equation}
\text{Comm}_{\text{ZeRO-3}} = \left( \frac{N-1}{N} \right) M + \left( \frac{N-1}{N} \right) M + \left( \frac{N-1}{N} \right) M = 3 \left( \frac{N-1}{N} \right) M
\end{equation}
}
Dividing by $\text{Comm}_{\text{DDP}}$:
{\boldmath
\begin{equation}
\frac{\text{Comm}_{\text{ZeRO-3}}}{\text{Comm}_{\text{DDP}}} = \frac{3 \left( \frac{N-1}{N} \right) M}{2 \left( \frac{N-1}{N} \right) M} = \frac{3}{2} = 1.5
\end{equation}
}
Thus, ZeRO-3 trades a modest $50\%$ increase in communication volume for an $\mathcal{O}(N)$ reduction in static memory.
\end{proof}

\section{Foundational Architectural Questions and Epistemic Meta-Questions}

\subsection{Architectural Question 1: Overlapping Communication with Prefetching}
\textbf{Question:} How does FSDP avoid waiting for All-Gather communication during forward and backward passes?\\
\textbf{Answer:} FSDP implements double-buffered asynchronous prefetching: while GPU tensor cores are computing forward activations for layer $l$, a separate CUDA communication stream executes an asynchronous All-Gather for layer $l+1$'s weights. Because compute time typically exceeds transfer time over fast interconnects (NVLink), communication is 100\% hidden behind computation.

\subsection{Architectural Question 2: Checkpoint Resharding}
\textbf{Question:} How are ZeRO-3 / FSDP checkpoints saved and resumed when changing cluster size?\\
\textbf{Answer:} In ZeRO-3, each GPU writes only its local $1/N$ shard in parallel alongside a global index manifest. When resuming on a different cluster with $M$ GPUs, distributed checkpointing engines (such as PyTorch DCP, #241) re-assemble the global state or reshard shards into $1/M$ partitions on-the-fly.

\section{Algorithmic Implementation Specification}

\begin{algorithm}[H]
\caption{ZeRO-FSDP Memory and Communication Accounting}
\begin{algorithmic}[1]
\Require Total parameters $P$, number of devices $N$, stage $\in \{1, 2, 3\}$, bytes per param $B$, optimizer
\Ensure Memory per device $\mathcal{M}_{\text{dev}}$, saving factor $\mathcal{S}$, sharded params $P_{\text{shard}}$, communication $C$
\State $B_{\text{opt}} \gets (\text{optimizer} == \text{"adam"}) \; ? \; 12.0 : 4.0$
\State $M_{\text{param}} \gets P \cdot B, \quad M_{\text{grad}} \gets P \cdot B, \quad M_{\text{opt}} \gets P \cdot B_{\text{opt}}$
\State $M_{\text{ddp}} \gets M_{\text{param}} + M_{\text{grad}} + M_{\text{opt}}$
\If{$\text{stage} == 1$}
    \State $M_{\text{dev}} \gets M_{\text{param}} + M_{\text{grad}} + (M_{\text{opt}} / N)$
    \State $C \gets 2.0 \cdot ((N - 1.0) / N) \cdot M_{\text{grad}}, \quad P_{\text{shard}} \gets P$
\ElsIf{$\text{stage} == 2$}
    \State $M_{\text{dev}} \gets M_{\text{param}} + (M_{\text{grad}} / N) + (M_{\text{opt}} / N)$
    \State $C \gets ((N - 1.0) / N) \cdot M_{\text{grad}}, \quad P_{\text{shard}} \gets P$
\Else
    \State $M_{\text{dev}} \gets (M_{\text{param}} / N) + (M_{\text{grad}} / N) + (M_{\text{opt}} / N)$
    \State $C \gets 3.0 \cdot ((N - 1.0) / N) \cdot M_{\text{param}}, \quad P_{\text{shard}} \gets \lceil P / N \rceil$
\EndIf
\State $\mathcal{M}_{\text{dev}} \gets M_{\text{dev}} / (1024^2)$
\State $\mathcal{S} \gets M_{\text{ddp}} / M_{\text{dev}}, \quad C_{\text{mb}} \gets C / (1024^2)$
\State \Return $(\mathcal{M}_{\text{dev}}, \mathcal{S}, P_{\text{shard}}, C_{\text{mb}})$
\end{algorithmic}
\end{algorithm}

\section{Big-$\Theta$ Complexity Analysis}

\begin{itemize}
    \item \textbf{Memory Calculation Time}: $\Theta(1)$ constant time operations.
    \item \textbf{Per-Device Parameter Footprint}: $\Theta(P / N)$ memory for ZeRO-3.
    \item \textbf{Per-Device Communication Volume}: $\Theta(P)$ bytes per step.
    \item \textbf{Space Savings}: Asymptotically optimal $\Theta(N)$ reduction.
\end{itemize}

\section{Single Canonical Repository Reference}

The complete deterministic scalar implementation, verification suite, and unit test benchmarks are published at:
\begin{center}
\url{https://github.com/Chief-Strategist-J/llm-obs-infra}
\end{center}

\end{document}
